% =====================================================================
%  Night Flash Edge -- research paper
%  Build:  tectonic -X compile paper.tex --outfmt pdf
%         (or ../latex/build.sh paper, from docs/latex/)
% =====================================================================
\documentclass[11pt,a4paper]{article}

\usepackage[margin=25mm]{geometry}
\usepackage{amsmath,amssymb,booktabs,graphicx,longtable,array,enumitem}
\usepackage[hidelinks]{hyperref}
\usepackage{fancyhdr,titlesec,microtype}
\usepackage{xcolor}
\usepackage{listings}

\definecolor{codebg}{RGB}{247,247,249}
\definecolor{rulegray}{RGB}{200,200,205}
\definecolor{accent}{RGB}{30,60,110}

\lstset{
  backgroundcolor=\color{codebg},
  basicstyle=\ttfamily\small,
  breaklines=true,
  frame=single,
  rulecolor=\color{rulegray},
  columns=fullflexible,
  keepspaces=true,
  showstringspaces=false,
}

\titleformat{\section}{\Large\bfseries\color{accent}}{\thesection}{0.7em}{}
\titleformat{\subsection}{\normalsize\bfseries}{\thesubsection}{0.6em}{}
\titlespacing*{\section}{0pt}{14pt}{6pt}
\titlespacing*{\subsection}{0pt}{10pt}{4pt}

\pagestyle{fancy}
\fancyhf{}
\fancyhead[L]{\small\color{gray}Night Flash Edge}
\fancyhead[R]{\small\color{gray}\thepage}
\renewcommand{\headrulewidth}{0.4pt}

\newcommand{\code}[1]{\texttt{#1}}
\newcommand{\flag}[1]{\texttt{--#1}}
\newcommand{\method}{\textsc{Night Flash Edge}}

\title{\vspace{-2.2cm}\textbf{Compact Unpaired Style Transfer for a}\\
\large Disposable-Camera Aesthetic, with Synthesised Source Domains\\[6pt]
\large and Held-Out Model Selection}
\author{}
\date{\today}

\begin{document}
\maketitle
\vspace{-1.1cm}

\begin{abstract}
\noindent
Learning a photographic aesthetic from a target domain alone is the binding
constraint of unpaired image-to-image translation, and it is usually left
unsolved. We present \method, a $4.02$ M-parameter generator that learns a
disposable-camera / Night Flash aesthetic for a native mobile camera, subject to
two hard requirements: real-time execution on device, and export to Core ML.

We make three contributions. First, we identify and correct a systematic bias in
the standard unpaired objective: when an $L_1$ content penalty is evaluated on a
\emph{stochastic} output, it penalises the stochastic texture the model is being
asked to produce. We measure the penalty at $1.9\times$ the magnitude of the
entire learned residual, and show that separating the $L_1$ and edge halves across
different images removes it. Second, we introduce a synthetic degradation pipeline
that manufactures a source domain from the target photographs by simulating a
non-flash camera, in two variants: one with exact correspondence, which admits
supervised reconstruction and real quality metrics, and one without, which
provably removes the inverse-mapping shortcut the first variant admits. Third, we
supply a validation regime that makes model selection possible at all, using
reproducible held-out strips, a per-epoch metric history, and best-checkpoint
retention --- none of which a loss curve provides.

We report measurements for every claim, including two places where measurement
falsified our initial design claims, and report the resulting corrections.
\end{abstract}

\tableofcontents
\newpage

% =====================================================================
\section{Introduction}

A mobile camera application typically ships a deterministic colour transform
plus a finishing chain, and can render a fixed aesthetic reliably. What it cannot
do is \emph{adapt}: choose how much flash falloff a given scene wants, how
strongly to lift a shadow, where grain should sit. A learned filter can, provided
it is small enough to run in real time and simple enough to export to a hardware
accelerator.

The obstacle is data. A photographic aesthetic is a property of lighting and film
response, not of content, so the pairs one would want --- the same scene with and
without the treatment --- do not exist. Training must therefore be unpaired, and
unpaired training brings a specific difficulty that this work addresses directly:
\emph{with only a target domain, there is no source domain to translate from}.

Three questions structure the paper. \S\ref{sec:loss} asks whether the standard
unpaired objective is correctly specified when the generator emits stochastic
texture, and finds that it is not. \S\ref{sec:degrade} asks how a source domain
can be manufactured when only a target exists. \S\ref{sec:val} asks how model
selection is possible without paired validation data, which is the precondition
for knowing whether any of the above worked.

\paragraph{Contributions.}
\begin{enumerate}[leftmargin=1.5em]
  \item A measurement and correction of texture bias in the unpaired content
        penalty (\S\ref{sec:loss}).
  \item A nine-stage camera-degradation model for synthesising a source domain,
        with a paired and an unpaired pairing scheme, and an analysis of what it
        can and cannot achieve (\S\ref{sec:degrade}).
  \item A held-out validation regime with deterministic reproduction,
        per-image comparability, a metric history, and best-checkpoint
        selection (\S\ref{sec:val}).
  \item A released, tested implementation in which all $28$ architecture
        $\times$ head combinations train, infer and convert to Core ML
        (\S\ref{sec:impl}).
\end{enumerate}

% =====================================================================
\section{Method}

\subsection{Model}

The generator predicts a bounded residual on the input:
\begin{equation}
  x_{\text{out}} = \operatorname{clamp}\!\left(x_{\text{in}} + \Delta_\theta(x_{\text{in}}), 0, 1\right),
  \qquad \lVert\Delta_\theta\rVert_\infty \le \rho .
  \label{eq:residual}
\end{equation}

Eq.~\eqref{eq:residual} is the load-bearing design decision. The only way to
produce a pixel is to perturb a pixel that already exists, so faces, text and
geometry are structurally protected rather than protected by a penalty term. It
also fixes the failure direction: an under-trained model degrades to a
near-identity filter, which is the correct behaviour for a camera filter.

The trunk is a convolutional encoder/decoder with local-window self-attention
(\code{EdgeViT}) and a single overlapping cross-attention stage (\code{HAT}) for
global context, at $4.02$ M parameters. Three photometric capabilities --- halation,
a subject gate, and per-region exposure --- are each gated by a
zero-initialised mix scalar, so a fresh model is exactly the plain residual and
each capability engages only once training motivates it.

The head determines the permitted parameterisation. We use a frequency-split head
by default, predicting low, mid and high bands separately, on the grounds that
grain and halation are high-frequency phenomena that a per-pixel head
represents poorly.

\subsection{Objective}

\begin{equation}
  \mathcal{L}_G = \mathcal{L}_{\mathrm{adv}} + \lambda_c \mathcal{L}_{\mathrm{content}}
  + \lambda_m \mathcal{L}_{\mathrm{moment}} + \lambda_i \mathcal{L}_{\mathrm{identity}}
  + \lambda_e \mathcal{L}_{\mathrm{effect}},
  \label{eq:obj}
\end{equation}
with $(\lambda_c, \lambda_m, \lambda_i, \lambda_e) = (1.5, 0.35, 0.75, 0.02)$ and a
PatchGAN discriminator at unit weight. The discriminator is the only term that
sees targets as targets; every other term is a constraint on the generator.

% =====================================================================
\section{A correction: texture bias in the content penalty}
\label{sec:loss}

\subsection{The observation}

The generator emits a deterministic residual, which a seeded procedural stage
then turns into a stochastic image by adding grain, leak, exposure, saturation
and vignette. Let $x_{\text{det}}$ denote the pre-effect output and
$x_{\text{styled}}$ the post-effect output.

The content term is a convex combination of an $L_1$ half and a Sobel edge half.
Both halves were originally evaluated on $x_{\text{styled}}$. Because $L_1$ is
texture sensitive, this charges the output for the stochastic texture.

\subsection{Measurement}

We train a model on a small pool and measure both penalties directly. The
predicate ``$L_1$ on the styled image'' and ``$L_1$ on the deterministic image''
differ by construction, but the interesting quantity is the magnitude:

\begin{table}[h]
\centering
\begin{tabular}{@{}lr@{}}
\toprule
Quantity & Value \\
\midrule
$\mathrm{content}(x_{\text{det}}, x_{\text{src}})$ & 0.01451 \\
$\mathrm{content}(x_{\text{styled}}, x_{\text{src}})$ & 0.04237 \\
\textbf{Penalty incurred by the effects} & \textbf{$+0.02786$} \\
\midrule
$\lVert x_{\text{det}} - x_{\text{src}} \rVert_1$ & 0.00891 \\
$\lVert x_{\text{styled}} - x_{\text{det}} \rVert_1$ & 0.01601 \\
\bottomrule
\end{tabular}
\caption{The stochastic effects are charged roughly $1.9\times$ more by the
content term than the entire learned residual is charged for the whole style
transfer.}
\label{tab:bias}
\end{table}

Table~\ref{tab:bias} is the whole argument. The adversarial term, at unit
weight, asks for grain. The content term, at weight $1.5$, penalises it, by an
amount nearly double to the residual's own contribution. The content term wins,
and the model learns to suppress the very effects its own effect head exists to
produce. This is the familiar failure where a squared or absolute error on a
stochastic target averages the stochastic component away.

\subsection{The correction}

The two halves belong to different comparisons. $L_1$ is a texture statistic and
must be evaluated where the texture is not stochastic; the Sobel edge response
is a structure statistic, is indifferent to fine texture, and should be
evaluated on the final image so that broken geometry is actually caught:

\begin{equation}
  \mathcal{L}_{\mathrm{content}}
    = 0.65\, L_1\!\left(x_{\text{det}},\, x_{\text{src}}\right)
    + 0.35\, \mathcal{E}\!\left(x_{\text{styled}},\, x_{\text{src}}\right).
  \label{eq:content}
\end{equation}

Eq.~\eqref{eq:content} leaves the intent of the original term --- stay close to
the source --- intact while removing the gradient that fights the texture.

\paragraph{Consequence for comparability.} Grain and halation become visibly
stronger, which is the intent, but any checkpoint trained under the previous
balance is no longer comparable. We therefore version the objective
(\code{LOSS\_VERSION}$=2$), record the version in every checkpoint and in
\code{config.json}, and warn on resume when it differs, so that the two halves of
a run are never silently pooled.

\subsection{Related failure avoided in the supervised term}

Where a correspondence exists, we add a reconstruction term. It is deliberately
evaluated on $x_{\text{det}}$ for the same reason: an $L_1$ term pointed at a
stochastic \emph{target} would average the target's own grain out of both sides.

% =====================================================================
\section{Synthesised source domains}
\label{sec:degrade}

\subsection{Motivation}

Unpaired translation assumes a source domain. When only a target domain exists,
one must be constructed, or the problem is not posed. The construction must be a
plausible capture of the same scenes by a different camera, so that the learned
map generalises to real photographs rather than to the construction's
fingerprints.

\subsection{Pipeline}

We simulate a non-flash camera by undoing a capture pipeline, resampling first so
that photometric operations act on clean pixels:

\begin{center}
\begin{tabular}{@{}cl@{}}
\toprule
\# & Stage \\
\midrule
1 & crop and small rotation, shared with the target \\
2 & downscale then upscale --- removes grain and fine detail \\
3 & Gaussian blur --- residual softening \\
4 & contrast compression about mid-grey \\
5 & black lift, $x \mapsto x + b(1-x)$, preserving white \\
6 & exposure gain \\
7 & white-balance shift, warmed \\
8 & saturation increase \\
9 & light sensor noise \\
\bottomrule
\end{tabular}
\end{center}

Stages 4--5 carry most of the shift. Compressing the tonal range is what removes
the flash falloff signature --- the bright subject on dark ground --- and it is
what makes the output read as an evenly-lit non-flash capture rather than merely
a dark one.

\subsection{Effect, measured}

Averaged over eight draws on a real photograph:

\begin{center}
\begin{tabular}{@{}lccc@{}}
\toprule
Statistic & \code{subtle} & \code{moderate} & \code{strong} \\
\midrule
high-frequency energy (relative) & 0.52 & 0.30 & 0.21 \\
tonal range $p_{99}-p_1$         & 0.75 & 0.57 & 0.49 \\
warmth $R-B$                     & 0.040 & 0.078 & 0.121 \\
\bottomrule
\end{tabular}
\end{center}

\subsection{Two claims we had to retract}

\paragraph{Resampling does not remove halation.} We initially documented the
downscale/upscale stage as destroying halation bloom. A targeted probe refuted
this. Halation is a wide, low-frequency phenomenon; a low-pass filter has nothing
to remove, and blurring a glow makes it \emph{brighter} by spreading the core
into the surrounding ring. On a probe with a specular point and a halo, the
ring-to-peak luminance ratio rises from $0.57$ to $0.86$ as blur increases from
$0$ to $4$. The pipeline attenuates halation, driven by tonal compression rather
than filtering, which suffices for our purpose --- the model must learn to add
glow where the target has it --- but would not suffice for constructing a source
with provably zero halation. That would require re-rendering the illumination.

\paragraph{Too-strong degradation is a defect.} A first \code{strong} preset used
blur to $2.5$ and black lift to $0.20$. The output looked foggy rather than like a
photograph, and its injected noise was conspicuous against the softened
background. Training on such inputs teaches a model to expect a distribution that
never occurs at inference, which is a worse outcome than a weaker domain shift.
We pulled the ranges back until the strongest profile still reads as a believable
overcast capture. Plausibility, not magnitude, is the design objective.

\subsection{Two pairing schemes}

\begin{center}
\begin{tabular}{@{}llll@{}}
\toprule
Scheme & Source & Target & Correspondence \\
\midrule
\code{paired}  & $\mathrm{deg}(x_i)$ & $x_i$ & exact, same file \\
\code{sameset} & $\mathrm{deg}(x_j),\ j \neq i$ & $x_i$ & none, one pool \\
\bottomrule
\end{tabular}
\end{center}

In both, the shared geometric warp is applied to both images first and the
degradation runs afterwards on the source only, so any crop, flip or rotation is
common to the pair and only the simulated camera differs.

\subsection{The inverse-mapping shortcut}

The \code{paired} scheme admits a shortcut. Because the source is a deterministic
function of the target, the network can reduce its loss by learning
$\mathrm{deg}^{-1}$ rather than the aesthetic. It will score well on
freshly-degraded inputs and fail on real photographs, and the failure is invisible
unless sought.

The \code{sameset} scheme removes the shortcut by construction: $x_i$ and $x_j$
are unrelated photographs, so no pointwise relationship exists to invert and the
model must learn the conditional distribution of the look. It also doubles the
usable pairings, since every image serves as both source and target.

The price is that without correspondence no supervised reconstruction is
available and no real quality number can be computed, so only distribution-level
metrics apply. We accept this asymmetry rather than pretending it away, and the
practical recommendation is to run both schemes as parallel arms and let the
evidence decide.

\subsection{Variation across images, epochs and runs}

Parameters are drawn per image; the dataset seeds its generator with the epoch
index, so the draw also varies per epoch at no cost; and the run seed shifts the
stream. Two further mechanisms are provided. An \emph{identity path}, taken with
probability $0.10$, feeds the target through undegraded, anchoring the identity
case and widening the input distribution. A \emph{curriculum}, on by default,
ramps the profile \code{subtle} $\rightarrow$ \code{moderate} $\rightarrow$
\code{strong}, giving early epochs an easy mapping and later epochs robustness.

% =====================================================================
\section{Held-out validation and model selection}
\label{sec:val}

\subsection{The problem}

A discriminator loss is not a quality signal, and neither is a generator loss.
Without a validation signal there is no basis for choosing between checkpoints, no
basis for detecting overfitting, and no basis for deciding whether a
hyperparameter change helped. In this setting the difficulty is sharper still,
because the data is unpaired by construction.

\subsection{Deterministic reproduction}

Two independent properties must be pinned, and the second is unconditional.

\emph{Which} images appear is set by \code{val\_select}; we default to
\code{fixed}, because per-image progress is only legible if the same images
recur, and because re-showing the same handful would otherwise hide a model that
has memorised them.

\emph{How} each image is degraded and seeded is always pinned, with seeds derived
from the filename by SHA-256 rather than from the epoch-dependent random stream.
This is not an optimisation. A validation image degraded with a fresh draw each
epoch is a different problem each time, and any apparent improvement is noise.
Python's built-in \code{hash()} is salted per interpreter and is therefore
unusable for this purpose; cross-session stability requires a cryptographic
digest.

The resulting artefact is one strip image per epoch: one image per row, cells
ordered input, output, target, so a given image's input and output are adjacent.
The header carries the epoch, mode, active profile and aggregate metrics.

\subsection{From measurement to selection}

The strip is for a person; the CSV history is for the log. One row per epoch of
PSNR, SSIM, $L_1$ and effect magnitude makes the curve plottable, and --- more
importantly --- makes checkpoint selection possible. Before this, only
\code{latest.pt} existed: a run peaking at epoch $40$ and decaying by epoch $60$
had no way to recover its best weights. We now retain a \code{best.pt} on
improvement and prune old per-epoch checkpoints, since $60$ epochs of a
$17$ MB checkpoint is a gigabyte that also slows every snapshot.

We write \code{best.pt} only in the \code{paired} scheme. In \code{sameset},
PSNR compares two unrelated photographs, so ranking on it would be meaningless,
and we decline to emit a checkpoint that appears authoritative but is not.

\subsection{A shared-warp defect}

The first implementation of the shared geometric warp sampled scale, rotation and
flips once but drew the crop \emph{offset} from the shared random stream inside
each application, so the source and target received different offsets. This is
harmless for unpaired data, where the images are unrelated anyway, and it
silently destroys the correspondence for paired data. We sample the offset once
as a fraction and clamp it per image, so differently sized images still receive
the same relative crop. We report it because the class of bug --- a
correspondence silently destroyed by a shared mutable stream --- is easy to
introduce and invisible without a test that asserts byte-identical outputs from
identical inputs.

% =====================================================================
\section{Deployment}

\subsection{Inference}

Whole-frame inference is the default: the image is letterboxed, styled, and
restored to its exact original dimensions, with no cropping and no stretching.
We verify a landmark round trip at $0.0017$ worst positional error with output
size identical to input.

Tiled inference, for native-resolution detail, crossfades overlapping tiles under
a raised-cosine window. The tile size must match the training crop; a mismatch is
a distribution shift, and the implementation warns rather than blocks. Measured
cost on a $4032\times3024$ source at $256$ px tiles is $336$ tiles and $22.0$ MP
at $25\%$ overlap, rising to $713$ tiles and $46.7$ MP at $50\%$. The relevant
quantity is the overlap \emph{ratio}, not its absolute value.

Frame borders do not fade, which is the usual failure of naive windowed blending.
The window carries a positive floor and edge tiles extend past the frame, so the
frame border is never a tile border; a $255$-valued border strip survives at
$254.6/255$ against a worst-case error of $1.0$.

\subsection{Core ML}

All $28$ architecture $\times$ head combinations convert, with tensor-I/O mean
absolute error $0.00030$. Several choices exist solely to satisfy the export
boundary and should not be refactored away: explicit attention matmuls in place
of the fused multi-head operator, a flat $(s^3, 3)$ LUT table, explicit trilinear
gathers in place of \code{grid\_sample}, and shape integers precomputed at
construction because \code{torch.jit.trace} turns \code{x.shape[1]} into a
zero-dimensional tensor that breaks pooling and \code{min}/\code{max}.

\subsection{Training/inference gap}

Synthesised source domains narrow the input distribution relative to arbitrary
user photographs. We mitigate this on the training side --- wide randomisation,
the identity path, and a strong profile --- and we explicitly reject
inference-time input normalisation, which risks double-degrading a user image.

% =====================================================================
\section{Implementation and verification}
\label{sec:impl}

$264$ tests pass. Beyond shape and conversion checks, the suite pins the
behaviours that would otherwise fail silently: the content-loss split, including
that the $L_1$ half is insensitive to the styled image and the edge half is not;
that a shared warp yields byte-identical source and target from identical inputs;
that validation degradation is invariant across epochs while training degradation
is not; that the \code{sameset} scheme never pairs an image with itself and
degrades gracefully on a single-image pool; that \code{best.pt} is not written
without a correspondence; and an abstract-syntax check that no module unpacks a
generator call without the architecture-agnostic helper, a defect that had already
shipped once in the evaluation entry point.

\section{Conclusion}

The most consequential finding here is unglamorous: a standard loss term, applied
in a plausible place, silently cancelled the component of the model it was
supposed to support. It was found only by measuring the penalty each stage
incurits rather than by inspecting the loss curve, which continued to fall. The
second is that a validation regime is a prerequisite for knowing whether anything
worked, and that without a correspondence it must be built from
distributional evidence instead of reconstruction metrics. Both are, we think,
general lessons for small unpaired systems rather than facts about this
particular aesthetic.

\paragraph{Limitations.} The degradation model attenuates rather than eliminates
halation, and the strongest profile is a judgement call about plausibility. The
\method{} aesthetic is narrow by construction; multiple distinct looks would
require separate checkpoints or a style embedding. Unpaired training cannot
guarantee identity preservation, and the held-out metrics reported here are
quality proxies rather than perceptual ground truth.

% =====================================================================
\appendix
\section{Reproduction}

\begin{lstlisting}
# train, paired scheme, strip and best.pt every epoch
nf train --target-dir POOL --name a --epochs 60 --val-count 8

# the same pool with the inverse-mapping shortcut removed
nf train --target-dir POOL --name b --epochs 60 --train-mode sameset

# stronger degradation, no curriculum, EMA for flicker
nf train --target-dir POOL --name c --degrade-profile strong \
  --no-degrade-curriculum --ema-decay 0.999

# per-stage bounds, overriding the profile
nf train --target-dir POOL --name d --degrade-contrast-min 0.3 \
  --degrade-contrast-max 0.4 --degrade-resize-min 0.4
\end{lstlisting}

\begin{lstlisting}
# whole frame, exact original size (default)
s.infer("photo.jpg")

# tiled; match tile to image_size
s.infer("photo.jpg", tile=256, tile_overlap=64)

# downscale the long edge, aspect preserved, then tile
s.infer("12mp.jpg", max_edge=1024, tile=256)
\end{lstlisting}

\end{document}
